\documentclass[11pt,a4paper]{article}
\usepackage{gtu-el-yazisi}
\begin{document}
\notbilgi{Elektronik II}{İçsel kazanç, kaskod yükselteç ve akım aynası}{Ders-2.pdf}
\textit{Not: Orijinal dosyadaki basılı sayfalar (s. 9) dijitalleştirilmedi; yalnızca el yazısı sayfalar aktarıldı.}

\sayfa{1}
\elcizimi[0.9]{sekiller/Ders-2-s1-1.png}[1]

\sekilnotu[1]{Bu sayfanın orta kısmı basılı kitap materyali içeriyor (aktif yüklü CS/CE yükselteç devreleri, küçük işaret modelleri ve (7.1)--(7.6) denklemleri); dijitalleştirilmedi.}

$1\,\mathrm{mA}$, $V_{OV}$, $V_T = 0{,}025$
\begin{align*}
A_{v\,(\text{mosfet})} &= g_m r_o = \frac{2I_D}{V_{OV}}\cdot\frac{V_A' L}{I_D} = \frac{2V_A' L}{V_{OV}}\\
A_{v\,(\text{BJT})} &= g_m r_o = \frac{I_C}{V_T}\cdot\frac{V_A}{I_C} = \frac{V_A}{V_T}
\end{align*}
\textcolor{kirmizikalem}{$\uparrow\,\uparrow$ mosfete göre yüksek}

\sayfa{2}
\[ I_D = \tfrac12\,\mu_n C_{ox}\,\frac{W}{L}\,(V_{GS}-V_{Tn})^2 \]
\sekilnotu[2]{Bu sayfa basılı kitap materyali içeriyor (7.2.2 The Intrinsic Gain, (7.7)--(7.15) denklemleri); dijitalleştirilmedi. Üzerindeki kırmızı kalem işaretleri: (7.14)'te $A_0$ yanında aşağı ok, $V_{OV}$ yanında yukarı ok; aşağıdaki $I_D$ denkleminde $I_D$ ve karenin üstünde yukarı ok, $L$'nin altında aşağı ok.}
\begin{equation*} g_m = \frac{I_D}{V_{OV}/2} \tag{7.10} \end{equation*}
\[ I_D = \tfrac12\,\mu_n C_{ox}\,\frac{W}{L}\,\bigl(\underbrace{V_{GS}-V_{Tn}}_{\textcolor{kirmizikalem}{V_{OV}}}\bigr)^2 \]
Gerilim kaynağının iç direnci sıfırdır.\\
Akım kaynağının iç direnci sonsuzdur.

\elcizimi[0.5]{sekiller/Ders-2-s2-1.png}[2]

\sayfa{3}
\elcizimi[0.8]{sekiller/Ders-2-s3-1.png}[3]



\sekilnotu[3]{Sayfanın ortası basılı kitap materyali içeriyor ((7.16)--(7.18) denklemleri, $Q_1$/$Q_2$'li aktif yüklü yükselteç devresi); dijitalleştirilmedi. Üzerindeki kırmızı kalem eklemeleri: $Q_2$ ve $V_{DD}$ kırmızı elipsle çevrelenmiş, $Q_2$ gate'inde ``$V_G$'', ``Active load'' yazısı; sağda çarpı ile iptal edilmiş küçük çizim ($V_{DD}$, $g_m v_{gs2}$ kaynağı, $r_{o2}$).}

\elcizimi[0.8]{sekiller/Ders-2-s3-2.png}[3]

\elcizimi[0.9]{sekiller/Ders-2-s3-3.png}[3]

\sayfa{4}
{\color{kirmizikalem}
\[ \frac{v_{out}}{v_{in}} = \frac{-g_{m1}\,v_{gs1}\,(r_{o1}\parallel r_{o2})}{v_{gs1}} = -g_m\,(r_{o1}\parallel r_{o2}) \]
\[ \text{if } r_{o1}=r_{o2}=r_o \qquad\qquad A_v = -\tfrac12\, g_m r_o \]
\[ A_v = \frac12\cdot\frac{2I_D}{V_{OV}}\cdot\frac{V_A' L_1}{I_D}\footnotemark \]
\footnotetext{Orijinalde $\tfrac12$'deki 2 ile $2I_D$'deki 2 ve pay/paydadaki $I_D$ çizgiyle sadeleştirilmiş.}
\[ \uparrow A_v = \frac{V_A' L_1\ \uparrow}{V_{OV1}\ \downarrow} \]
\begin{center}
$\swarrow$\\
$V_{DS} \le V_{GS} - V_{TH}\ \nearrow$ sabit\\
$\downarrow$\\
fine tuning
\end{center}
}
\[ I_D = \tfrac12\,\mu_n C_{ox}\,\frac{\uparrow W}{\uparrow L}\,(V_{OV})^2 \]
$I_D$ azalmaması için

Onun yerine multistage yapalım

\sekilnotu[4]{Sayfanın alt kısmı basılı kitap materyali içeriyor (Chapter 7, Building Blocks of Integrated-Circuit Amplifiers; $Q_1$'li devre ve $g_{m1}v_i \parallel r_{o1}$ eşdeğeri (a), kaskat bağlı $Kr_{o1}$'li devre (b)); dijitalleştirilmedi.}

\sayfa{5}
\section{The Cascode Amplifier}
\subsection{Cascoding}

\sekilnotu[5]{Basılı kaskod şeması ($Q_2$ üstte, $Q_1$ altta, $V_{G2}$, $v_i$, ``To Load'', $Kr_{o1}$, $r_{o1}$ okları) dijitalleştirilmedi; üzerindeki el eklemeleri: en üstte $V_{DD}$ (yukarı ok) ucuna bağlı $I_{bias}$ akım kaynağı (daire, ok aşağı) ve yanında paralel direnç, yazı $R_{CB}=\infty$; bu kaynak $Q_2$'nin drain ucuna bağlanmış. $Q_1$ gate'i ``G'', $Q_1$ drain düğümü ($Q_2$ source) ``D'', $Q_1$ source (toprak) ``S'' olarak etiketlenmiş.}

\textcolor{kirmizikalem}{?? ödev}

\elcizimi[0.9]{sekiller/Ders-2-s5-1.png}[5]
\[ \textcolor{kirmizikalem}{V_x = -v_{gs2}} \]

\elcizimi[0.9]{sekiller/Ders-2-s5-2.png}[5]

\textcolor{kirmizikalem}{\underline{ödev}}

\sayfa{7}
\section{7.4 IC biasing}
\subsection{Current mirror}

\elcizimi[0.65]{sekiller/Ders-2-s7-1.png}[7]

{\color{turuncukalem}
\[ \frac{I_{D1} = I_{ref} = \tfrac12\mu_n C_{ox}\frac{W}{L}(V_{GS1}-V_{TH1})^2(1+\lambda V_{DS1})}{I_{D2} = I_O = \tfrac12\mu_n C_{ox}\frac{W}{L}(V_{GS2}-V_{TH2})^2(1+\lambda V_{DS2})}\footnotemark \]
\footnotetext{Orijinalde pay ve paydadaki ortak terimler ($\tfrac12\mu_n C_{ox}\frac{W}{L}$ ve $V_{TH}$) çizgiyle sadeleştirilmiş.}
}
\[ \frac{I_{ref}}{I_{D2}} = 1 \qquad\qquad I_{ref} = I_{D2} \]

What happen when we put a load ??

\elcizimi[0.65]{sekiller/Ders-2-s7-2.png}[7]

Bu ihtimalde akımlar aynı ve

\sayfa{8}
Yükler aynı ise sadece o gerilimler aynıdır. $V_{DS1}=V_{DS2}$

\elcizimi[0.9]{sekiller/Ders-2-s8-1.png}[8]

\elcizimi[0.9]{sekiller/Ders-2-s8-2.png}[8]



\elcizimi[0.65]{sekiller/Ders-2-s8-3.png}[8]

\sayfa{10}
\elcizimi[0.9]{sekiller/Ders-2-s10-1.png}[10]

\elcizimi[0.85]{sekiller/Ders-2-s10-2.png}[10]

\end{document}
